\documentclass[11pt]{article}

\usepackage[margin=1in]{geometry}
\usepackage[T1]{fontenc}
\usepackage[utf8]{inputenc}
\usepackage{lmodern}
\usepackage{microtype}
\usepackage{amsmath,amssymb,bm}
\usepackage{graphicx}
\usepackage{booktabs}
\usepackage{array}
\usepackage{tabularx}
\usepackage{caption}
\usepackage{subcaption}
\usepackage{enumitem}
\usepackage[hidelinks]{hyperref}
\usepackage{float}

\setlist[itemize]{leftmargin=1.4em}
\setlist[enumerate]{leftmargin=1.6em}
\captionsetup{font=small,labelfont=bf}
\renewcommand{\arraystretch}{1.18}

\title{\textbf{The Adaptive Boundary Equation:}\\
A Minimal Active-Boundary Model for Self-Organizing Information--Entropy Dynamics}
\author{Andrew Poinçot\\Cripsis Logic}
\date{Revised draft v10.9.2: Published May 8, 2026\\Suggested arXiv category: nlin.AO; secondary: q-bio.NC}

\begin{document}
\maketitle

\begin{abstract}
Adaptive systems persist by regulating information-bearing activity and entropy or complexity across boundaries. This paper proposes the Adaptive Boundary Equation (ABE), a normalized two-variable phenomenological model
\begin{equation*}
\frac{dI}{dt}=\gamma S-\lambda I-\beta I^3,\qquad
\frac{dS}{dt}=\eta I-\xi S,
\end{equation*}
with positive parameters. The cubic saturation term is the minimal nonlinearity required to bound the active state. The model yields a dimensionless viability number
\begin{equation*}
R=\frac{\gamma\eta}{\xi\lambda}.
\end{equation*}
The quiescent boundary is stable for $R<1$. For $R>1$, the quiescent state loses stability and a stable positive active branch appears:
\begin{equation*}
I^*=\sqrt{\frac{\gamma\eta/\xi-\lambda}{\beta}},\qquad S^*=\frac{\eta}{\xi}I^*.
\end{equation*}
Linear stability analysis and numerical verification confirm the threshold, active-branch stability, amplitude relation, and critical-slowing behavior near the viability transition. A worked synthetic operational demonstration in the style of an adaptive computational governor reproduces the predicted threshold transition and amplitude scaling. The contribution is not a new dynamical-system normal form, but a compact adaptive-boundary interpretation of familiar nonlinear machinery, together with explicit falsification criteria and one synthetic operational demonstration.
\end{abstract}

\noindent\textbf{Keywords:} adaptive systems; active boundaries; information dynamics; entropy production; self-organization; viability threshold; critical transitions; point processes; nonlinear dynamics; phenomenological models; active inference

\clearpage
\section{Introduction}
The central problem of adaptive systems is not merely how order appears, but how order persists under flux~\cite{schrodinger1944life}. Biological cells, neural circuits, institutions, markets, and computational governors remain viable only by regulating what crosses their boundaries, what is retained, what is dissipated, and what is forgotten. Each domain has its own microscopic mechanisms. Yet at the phenomenological level, many such systems exhibit similar signatures: threshold behavior, bounded active states, relaxation after perturbation, and sometimes damped oscillatory recovery.

Existing frameworks capture parts of this problem. Dissipative-structure theory~\cite{prigogine1967thermo,nicolis1977self} explains the emergence of order far from equilibrium. Maximum-entropy-production arguments~\cite{dewar2003info,martyushev2006mep} emphasize irreversible throughput. Free-energy formulations~\cite{friston2010free,buckley2017free} frame adaptation as inference under uncertainty. Control theory studies stabilization under feedback. Nonlinear dynamics supplies a mature language for thresholds, saturation, oscillation, and bifurcation.

The purpose of this paper is deliberately narrow. It does not claim a microscopic derivation of biological, economic, gravitational, or computational adaptation. It proposes a minimal active-boundary equation whose variables are interpreted as information-bearing activity and internal entropy or complexity. The paper then derives a viability threshold and shows what signatures would have to appear if the model is applicable to a given system.

This discipline is important. A model of this kind becomes scientifically useful only if it can be wrong. The ABE is therefore presented as a compact phenomenological target: simulate it, fit it, test whether the predicted threshold and amplitude scaling appear, and reject it where they do not.

\section{Variables and physical interpretation}
The model tracks two state variables.
\begin{itemize}
  \item $I(t)$: \textit{active information amplitude}. This is the boundary's organized information-bearing state, measured in a domain-specific unit such as effective memory strength, encoded state density, policy-bearing information, or normalized adaptive activity.
  \item $S(t)$: \textit{internal entropy or complexity}. This is the boundary's internal microstate diversity, disorder, adaptive complexity, or pressure generated by information processing and dissipated by regulation, compression, thermalization, or forgetting.
\end{itemize}

The normalized model uses five positive parameters.
\begin{itemize}
  \item $\gamma$: coupling gain from internal entropy or complexity into information-bearing activity.
  \item $\lambda$: regulatory decay of active information.
  \item $\beta$: nonlinear saturation coefficient limiting runaway activity.
  \item $\eta$: entropy or complexity generation rate per unit active information.
  \item $\xi$: entropy or complexity dissipation rate.
\end{itemize}

An optional forcing term $\Phi(t)$ represents external boundary drive. For threshold analysis $\Phi(t)$ is set to zero so that the question is endogenous viability: can the boundary maintain an active information-bearing state without continuous external forcing?

\section{The normalized Adaptive Boundary Equation}
The normalized active-boundary form of the ABE is
\begin{align}
\frac{dI}{dt}&=\gamma S-\lambda I-\beta I^3+\Phi(t), \label{eq:I}\\
\frac{dS}{dt}&=\eta I-\xi S. \label{eq:S}
\end{align}
Equation~\eqref{eq:I} says that information-bearing boundary activity grows when internal entropy or complexity provides useful adaptive drive, decays under regulatory loss, and saturates when activity becomes too large to remain linear. Equation~\eqref{eq:S} says that information-bearing activity generates entropy or complexity, while entropy or complexity dissipates at rate $\xi$.

The normalization removes an unnecessary response coefficient that would otherwise create dimensional ambiguity without changing the threshold. Time scaling or gain normalization can restore a domain-specific response rate if required. In this paper, the leading coefficient of the $I$ equation is absorbed into the definitions of $\gamma$, $\lambda$, $\beta$, and $\Phi$.

\section{Viability threshold}
Set $\Phi(t)=0$ and solve for equilibria. From Eq.~\eqref{eq:S},
\begin{equation}
S^*=\frac{\eta}{\xi}I^* . \label{eq:Sstar}
\end{equation}
Substituting Eq.~\eqref{eq:Sstar} into Eq.~\eqref{eq:I} gives
\begin{equation}
0=\left(\frac{\gamma\eta}{\xi}-\lambda\right)I^*-\beta (I^*)^3 . \label{eq:eqcond}
\end{equation}
The equilibria are
\begin{align}
I^*&=0,\qquad S^*=0, \label{eq:zero}\\
I^*&=\pm\sqrt{\frac{\gamma\eta/\xi-\lambda}{\beta}},\qquad S^*=\frac{\eta}{\xi}I^*. \label{eq:active}
\end{align}
The physically relevant positive active branch exists only when $\gamma\eta/\xi-\lambda>0$, equivalently
\begin{equation}
\eta>\frac{\xi\lambda}{\gamma}. \label{eq:threshold}
\end{equation}
The negative branch is the signed counterpart of the normal form; in applications where $I$ represents a nonnegative magnitude, the physical state is interpreted through $|I|$ or restricted to the positive branch. Define the dimensionless viability number
\begin{equation}
R=\frac{\gamma\eta}{\xi\lambda}. \label{eq:R}
\end{equation}
Then $R>1$ is the self-maintaining regime. The threshold is not imposed by interpretation; it follows directly from the equilibrium condition once the minimal saturation term is included. Below threshold, the only unforced equilibrium is quiescence. Above threshold, a bounded active branch becomes available.

\section{Linear stability analysis}
\subsection{Quiescent branch}
Linearizing Eqs.~\eqref{eq:I}--\eqref{eq:S} at $(I,S)=(0,0)$ with $\Phi(t)=0$ gives
\begin{equation}
J_0=\begin{pmatrix}-\lambda & \gamma \\ \eta & -\xi\end{pmatrix}. \label{eq:J0}
\end{equation}
The trace and determinant are
\begin{align}
\operatorname{Tr}(J_0)&=-(\lambda+\xi)<0,\\
\det(J_0)&=\lambda\xi-\gamma\eta=\lambda\xi(1-R). \label{eq:detJ0}
\end{align}
For $R<1$, the determinant is positive and the trace is negative, so the quiescent state is stable. For $R>1$, the determinant is negative, so the quiescent state is a saddle. In that regime, the quiescent boundary cannot be the stable self-maintaining state.

\subsection{Active branch}\label{sec:active-branch}
On the positive active branch, $\beta(I^*)^2=\gamma\eta/\xi-\lambda$. The Jacobian is
\begin{equation}
J^*=\begin{pmatrix}2\lambda-3\gamma\eta/\xi & \gamma\\ \eta & -\xi\end{pmatrix}. \label{eq:Jstar}
\end{equation}
The active-branch trace and determinant are
\begin{align}
T^*&=\operatorname{Tr}(J^*)=2\lambda-3\gamma\eta/\xi-\xi, \label{eq:Tstar}\\
D^*&=\det(J^*)=2(\gamma\eta-\lambda\xi), \label{eq:Dstar}\\
\eta_c&=\xi\lambda/\gamma. \label{eq:etac}
\end{align}
When $R>1$, $D^*>0$. The trace is also negative: since $R>1$ implies $\gamma\eta/\xi>\lambda$, we have $3\gamma\eta/\xi>3\lambda$, so
\begin{equation}
T^*<2\lambda-3\lambda-\xi=-\lambda-\xi<0. \label{eq:Tneg}
\end{equation}
Thus the positive active branch is stable whenever it exists. The eigenvalues are
\begin{equation}
\mu_{1,2}=\frac{T^*\pm\sqrt{(T^*)^2-4D^*}}{2}. \label{eq:eigs}
\end{equation}
For positive parameters in the minimal ABE, the discriminant is nonnegative on the active branch. Thus the minimal model predicts stable real relaxation rather than generic underdamped oscillation. This is not a weakness; it is a constraint. If an empirical system shows persistent oscillatory recovery, the minimal ABE is incomplete for that system and must be extended with delayed feedback, inhibitory coupling, or an additional inertial or memory state.

The robust dynamic claim is therefore stability and relaxation to the active branch, not a universal oscillation law. The robust scaling claim is the active-state amplitude law in Eq.~\eqref{eq:active}.

A quantitative consequence of Eqs.~\eqref{eq:Tstar}--\eqref{eq:Dstar} is critical slowing as the threshold is approached from above. The active-branch determinant satisfies 
\begin{equation}
D^*=2\xi\lambda(R-1),
\end{equation}
\noindent so $D^*\to 0$ as $R\to 1^+$. Since the smaller-magnitude eigenvalue scales as $D^*/|T^*|$ near threshold, the dominant relaxation rate vanishes there: small perturbations decay arbitrarily slowly as the system approaches the viability transition. Absence of this critical-slowing signature in a calibrated system would be inconsistent with a stable active branch born at the predicted threshold.

\section{Dimensionless interpretation}
The model can be written in dimensionless form to expose its control logic. After rescaling time and amplitudes, the unforced dynamics can be represented qualitatively as
\begin{align}
\frac{di}{d\tau}&=a(s-i-i^3),\\
\frac{ds}{d\tau}&=b(Ri-s),
\end{align}
where $a$ and $b$ are positive dimensionless response rates. In this form, $R$ is the main control number. $R<1$ means dissipative and regulatory losses dominate entropy-supported activity. $R>1$ means adaptive drive is strong enough to support a bounded active boundary.

The analogy to reproduction numbers in epidemiology is useful only at the level of threshold logic: a dimensionless control number crosses one, changing the qualitative behavior of the system. The ABE does not import the microscopic assumptions of epidemic models.

\section{Falsifiable signatures}
The ABE makes three classes of falsifiable predictions. These are model signatures, not domain-specific numerical forecasts until the parameters are calibrated in a given system.
\begin{enumerate}
  \item \textbf{Threshold signature.} Systems should transition from externally supported activity to self-maintaining activity as $R$ crosses one. The variables $\gamma$, $\eta$, $\xi$, and $\lambda$ must be estimated from the target system or constrained by independent measurements.
  \item \textbf{Amplitude signature.} Above threshold, the active information amplitude should obey $I^*\propto (R-1)^{1/2}$ when cubic saturation is the leading nonlinearity. Equivalently, $(I^*)^2$ should be linear in $\eta-\eta_c$ for fixed $\gamma$, $\xi$, $\lambda$, and $\beta$.
  \item \textbf{Relaxation signature.} Under the minimal positive-coupling model, calibrated systems should relax through stable real modes toward the appropriate branch. Strong underdamped oscillation would not confirm the minimal ABE; it would indicate the need for a delayed, inhibitory, or higher-dimensional extension.
\end{enumerate}
These are simple enough to be wrong. That is a strength. The ABE is useful only if its threshold and scaling laws survive contact with measured or simulated systems.

\section{Candidate calibration domains}
The most immediate calibration domain is adaptive computational governance, where queue pressure, latency drift, rejection rates, routing decisions, and policy outcomes can be logged and replayed. In that setting, the variable $I$ may represent effective policy-bearing information or organized control state, $S$ may represent operational complexity or pressure, and $\Phi(t)$ may represent incoming workload. A worked synthetic example in this style is given in Sec.~\ref{sec:synthetic}.

Neural, biological, economic, and organizational systems may also provide calibration targets, but those applications require separate parameter estimation and should not be treated as solved by the present paper. A neural application would need defensible proxies for stored synaptic information
\clearpage
\noindent and firing-pattern entropy. An organizational application would need proxies for institutional knowledge, turnover, procedural complexity, and innovation or learning rate. Without such calibration, the ABE remains a hypothesis-generating model, not a domain-specific theory. This paper therefore treats cross-domain use as a research program, not as proof of universal applicability.

\section{Stochastic pulse realization}
The deterministic ABE can be interpreted as a mean-field law. A stochastic implementation can realize the same qualitative behavior through event pulses and adaptive clustering. The point-process view is useful for computational systems because it allows a boundary law to be tested with replayable synthetic event streams before being fit to real telemetry.

One natural realization uses three coupled model families.
\begin{itemize}
  \item \textbf{Poisson pulse emitters:} independent or conditionally independent event sources produce raw stochastic pulse timing.
  \item \textbf{Hawkes-style clustering:} pulse arrivals alter future pulse intensity, producing burst structure and event clustering rather than flat randomness~\cite{hawkes1971spectra}.
  \item \textbf{Markov damping:} a finite-state process regulates transitions between calm, active, burst, recovery, and saturated regimes, preventing uncontrolled excitation~\cite{norris1997markov,gardiner2009stochastic}.
\end{itemize}
In short: Poisson emitters provide raw pulse timing; Hawkes-style dynamics shape burst structure; Markov damping stabilizes transitions between burst cycles. The deterministic ABE is then interpreted as the coarse-grained boundary behavior produced by many such microscopic pulse events.

A full derivation of the deterministic ABE as a mean-field limit of a Poisson--Hawkes--Markov construction is left to future work. The bridge is plausible because mean-intensity equations for self-exciting point processes are well established~\cite{bremaud1996stability,daley2003point}, but it is not claimed as completed here.

\clearpage
\section{Numerical verification of analytic claims}
A minimal validation suite was run against the normalized equations. The goal of this section is not to validate any empirical domain. The goal is to verify that the simulations, algebra, and stability claims all use the same model. The checks evaluate the equilibrium branch, stability conditions, amplitude scaling, active-branch stability relation, and critical-slowing behavior near the viability transition. In all cases, the results match the analytical predictions.

\begin{figure}[H]
\centering
\includegraphics[width=1.0\textwidth]{figures/figure1_phase_portrait.pdf}
\caption{Phase-portrait comparison. Below threshold ($R<1$, left), trajectories relax to the quiescent branch at the origin. Above threshold ($R>1$, right), trajectories relax to one of the two active equilibria, while the origin becomes an unstable saddle. Equilibria are explicitly marked. This figure illustrates the qualitative state-space transition rather than empirical calibration.}
\label{fig:phase}
\end{figure}

A fourth numerical check verifies the critical-slowing prediction in Sec.~\ref{sec:active-branch}. For fixed
$\gamma=\xi=\lambda=\beta=1$, we swept $R$ above threshold and computed the
dominant relaxation time from the smallest-magnitude real eigenvalue of the active-branch
Jacobian. As predicted, the relaxation rate tends to zero as $R\to 1^+$, and the
corresponding relaxation time increases sharply near the viability threshold. The result is shown in Figure~\ref{fig:critical-slowing}.
These figures do not replace empirical validation. They remove a more basic failure mode: claiming a threshold or oscillation law that the stated equations do not actually produce.

\begin{figure}[H]
\centering
\includegraphics[width=0.85\textwidth]{figures/figure2_amplitude_scaling.pdf}
\caption{Active-branch amplitude scaling from ODE integration. For fixed $\gamma=\xi=\lambda=\beta=1$, numerical trajectories were integrated from random small initial conditions and allowed to settle. The measured endpoint obeys $(I_{\mathrm{end}})^2\approx\eta-\eta_c$, matching the analytic prediction.}
\label{fig:amplitude}
\end{figure}

\vspace{-1.0em}

\begin{figure}[H]
\centering
\includegraphics[width=0.85\textwidth]{figures/figure3_stability_check.pdf}
\caption{Stability check. Random positive parameter sets were sampled. For $R<1$, the quiescent-branch Jacobian has negative largest real eigenvalue. For $R>1$, the active-branch Jacobian has negative largest real eigenvalue, confirming the stability proof across the parameter space.}
\label{fig:stability}
\end{figure}

\clearpage
\section{Synthetic operational demonstration}\label{sec:synthetic}
This section presents a worked demonstration in the style of an adaptive computational governor, for instance, an admission-control or backpressure system in front of a service that processes incoming work under contention. The demonstration is synthetic: the underlying generator is the ABE itself with operational measurement noise added to the equilibrium states. The point is not to validate the ABE against a real system. The point is to show what the predicted signatures look like when extracted from telemetry-shaped data, and how a practitioner would test them.

\subsection{Operational mapping}
In the computational-governor setting, the variables and parameters take a concrete operational meaning.
\begin{itemize}
  \item $I(t)$: organized control state, the effective amount of policy-bearing information actively shaping admission, routing, or backpressure decisions. Practical proxies include the magnitude of the active control signal, the rate of policy-driven decisions, or the measured strength of the governor's response.
  \item $S(t)$: operational complexity or pressure, the current backlog of work, queue depth, or accumulated decision uncertainty. Practical proxies include normalized queue pressure, in-flight work, or rolling latency variance.
  \item $\gamma$: how strongly accumulated complexity drives the governor's control state. Tunable in many systems via control-loop gain.
  \item $\lambda$: how quickly the governor's control state decays in the absence of pressure. Tunable via decay schedules and policy half-lives.
  \item $\eta$: how much complexity is generated per unit of active control. Determined by workload mix and system architecture.
  \item $\xi$: how quickly accumulated complexity is dissipated by regulation, retries, or completion. Determined by downstream capacity and cleanup paths.
  \item $\beta$: the saturation strength preventing runaway control activity. Set by hard limits, ceilings, or rate caps.
\end{itemize}
In this mapping, the viability number $R=\gamma\eta/(\xi\lambda)$ is a dimensionless statement about whether the governor's feedback gain dominates its regulatory dissipation. $R<1$ corresponds to a system in which control state cannot self-maintain, any active control collapses without sustained external workload. $R>1$ corresponds to a system in which a bounded active control state can persist endogenously.

\subsection{Demonstration protocol}
The protocol is straightforward and applicable to any system that exposes the operational variables above:
\begin{enumerate}
  \item Estimate $\gamma$, $\lambda$, $\beta$, and $\xi$ from system configuration and short-time response measurements, such as control-loop gains, decay schedules, and dissipation rates.
  \item Sweep $\eta$ by varying the workload-to-control coupling, for example, changing admission aggressiveness, batch composition, or operator-set policy weights, while holding the other parameters fixed.
  \item For each sweep point, allow the system to settle and record the equilibrium state $|I|$.
  \item Plot $|I|$ against $R=\gamma\eta/(\xi\lambda)$ and $|I|^2$ against $\eta-\eta_c$. The threshold should appear at $R=1$, and the amplitude-squared relation should be linear above threshold.
\end{enumerate}

\subsection{Synthetic results}
Figure~\ref{fig:synthetic} shows the result of this protocol applied to data generated by integrating the ABE with $\gamma=\lambda=\beta=\xi=1$, sweeping $\eta$ across the threshold, and adding Gaussian measurement noise of standard deviation $0.04$ to each equilibrium state to emulate real-world telemetry. Sixty parameter sweep points are shown.

Two observations are worth highlighting. First, the threshold transition at $R=1$ is visible directly in the data without parameter fitting; it is a structural feature, not a fitted result. Second, the linear scaling of $|I|^2$ in $\eta-\eta_c$ above threshold provides a quantitative test that goes beyond the threshold itself: a system that crosses the threshold but does not show the predicted amplitude law would falsify the cubic-saturation form even while supporting the threshold structure.

This demonstration is synthetic and is therefore not evidence that any particular real system is well described by the ABE. It is evidence that the analytic predictions are observable in noisy operational-style data, and it provides a concrete protocol that other practitioners can apply to their systems.

\begin{figure}[H]
\centering
\includegraphics[width=1.0\textwidth]{figures/figure4_synthetic_demo.pdf}
\caption{Synthetic operational demonstration. Left: measured equilibrium $|I|$ as a function of $R=\gamma\eta/(\xi\lambda)$. Activity is consistent with zero below $R=1$ within the imposed measurement noise ($\sigma=0.04$), and rises sharply at the predicted threshold. Right: above-threshold equilibrium amplitude squared, plotted against $\eta-\eta_c$. The data fall along the predicted theoretical line, with residuals consistent with the imposed measurement noise.}
\label{fig:synthetic}
\end{figure}

\begin{figure}[H]
\centering
\includegraphics[width=1.0\textwidth]{figures/figure5_critical_slowing.pdf}
\caption{Critical slowing near the active-boundary threshold. For fixed $\gamma=\xi=\lambda=\beta=1$, the dominant relaxation time is computed from the smallest-magnitude real eigenvalue of the active-branch Jacobian as $R$ approaches $1$ from above. The relaxation time increases sharply near threshold, matching the analytical prediction that the dominant relaxation rate vanishes as $D^*=2\xi\lambda(R-1)$ tends to zero.}
\label{fig:critical-slowing}
\end{figure}

\section{Relationship to existing nonlinear systems}
The system in Eqs.~\eqref{eq:I}--\eqref{eq:S} shares its functional form with familiar two-variable nonlinear models, including FitzHugh--Nagumo-type relaxation systems~\cite{fitzhugh1961impulses,nagumo1962active}, cubic-saturation feedback models, and normal-form treatments near dynamical instabilities~\cite{cross1993pattern,kuramoto1984chemical}. In the signed-amplitude representation, the threshold has the structure of a supercritical pitchfork-type active-branch transition (see, e.g., \cite{strogatz2018nonlinear} for the canonical treatment); the novelty claimed here is the active-boundary interpretation, not the normal form itself. A reader trained in nonlinear dynamics should not read the ABE as a wholly new mathematical structure.

The contribution of this paper is instead interpretive and organizational. It assigns the variables and parameters a specific active-boundary meaning: information-bearing activity, entropy or complexity, regulatory decay, entropy generation, dissipation, and saturation. It then identifies $R=\gamma\eta/(\xi\lambda)$ as a compact viability number for this interpretation, gives falsifiable threshold, amplitude, and damping signatures, and demonstrates the predictions on synthetic operational data.

Three adjacent research traditions deserve explicit mention. First, Friston and collaborators have developed a body of work on Markov blankets and active inference~\cite{friston2013life,kirchhoff2018markov} that addresses the same broad question of how systems maintain themselves through boundaries that separate internal from external states. Their framework operates through variational free-energy minimization on a generative model rather than through phenomenological ODEs, and produces a different mathematical apparatus for a structurally similar question. Second, viability theory~\cite{aubin1991viability} formalizes self-maintenance through set-valued dynamics and viability kernels, providing a control-theoretic framework that overlaps with the present terminology of ``viability'' but employs different machinery. Third, work on phenomenological model reduction in systems biology~\cite{transtrum2015perspective,daniels2015automated} has shown that complex adaptive networks can often be characterized by a small number of dimensionless ratios derived from manifold-boundary approximations of microscopic kinetics, with adaptive systems sometimes reducing to a single timescale ratio. The present work differs from each of these traditions in mathematical approach: a minimal two-variable ODE with explicit threshold and amplitude predictions, falsifiable through three operational signatures, and calibrated through phenomenological parameter estimation rather than through variational inference, set-valued analysis, or microscopic reduction. The interpretive contribution is the active-boundary framing of an existing nonlinear normal form together with a viability number compact enough to be tested directly against operational data. This positioning is intentionally modest. The paper matters only if the interpretation helps connect measurable adaptive-boundary behavior to an analytically tractable model.

\clearpage
\section{Falsification criteria}
Table~\ref{tab:falsification} makes the falsification structure explicit. Each row gives a class of empirical observation that, if found in a calibrated system, would either falsify the minimal ABE or indicate the type of extension required. This table is not exhaustive, but it covers the central failure modes a reviewer or practitioner would test for.

\begin{table}[H]
\centering
\small
\caption{Falsification criteria for the minimal ABE. Each observation either falsifies the minimal model or indicates the specific extension required.}
\label{tab:falsification}
\begin{tabularx}{\textwidth}{>{\raggedright\arraybackslash}p{0.28\textwidth} >{\raggedright\arraybackslash}p{0.34\textwidth} >{\raggedright\arraybackslash}X}
\toprule
\textbf{Observation} & \textbf{Implication for the minimal ABE} & \textbf{Candidate response}\\
\midrule
Sustained nonzero activity in the unforced regime with $R$ clearly $<1$ & Falsifies the minimal model. Activity must be either externally forced, driven by a hidden state, or supported by a different threshold law. & Re-examine forcing assumptions; consider an extended ABE with bistability or hidden state.\\
\addlinespace
Quiescence persists despite $R$ well above $1$ & Falsifies the minimal model in that domain. The active branch should exist and be reachable. & Investigate basin-of-attraction issues, frozen initial conditions, or missing coupling.\\
\addlinespace
Active amplitude does not scale as $(I^*)^2\propto \eta-\eta_c$ near threshold & The leading nonlinearity is not cubic. Threshold structure may still hold; the saturation law is wrong. & Replace cubic saturation with the appropriate higher-order or bounded form.\\
\addlinespace
Persistent underdamped oscillation around the active branch & Falsifies the minimal positive-coupling ABE for that domain. The discriminant is forced negative by structure not in the model. & Extend with delayed feedback, inhibitory coupling, or an additional inertial state.\\
\addlinespace
No critical slowing as $R$ approaches $1$ from above & Inconsistent with a stable active branch born at threshold. Either $R$ is mis-estimated or the transition is not the one modeled. & Re-fit parameters; consider whether the empirical transition is a different bifurcation type.\\
\addlinespace
Threshold appears at a value of $R$ clearly different from $1$ & The chosen variable interpretation or normalization is wrong, even if the qualitative structure is right. & Revisit the operational mapping of $I$, $S$, $\gamma$, $\eta$, $\lambda$, and $\xi$ in the target system.\\
\bottomrule
\end{tabularx}
\end{table}

These criteria are deliberately operational. A system that survives all of them is consistent with the minimal ABE in the regime tested; a system that fails one of them tells the practitioner specifically how the model needs to be modified.

\clearpage
\section{Limitations}
\begin{enumerate}
  \item \textbf{Phenomenological status.} The ABE is not yet a microscopic derivation. Its parameters must be fit or derived for each domain.
  \item \textbf{Saturation choice.} Cubic saturation is minimal and analytically convenient. Other saturation laws may change amplitudes and prefactors while preserving the threshold structure.
  \item \textbf{Oscillation limits.} The minimal positive-coupling ABE does not produce a generic underdamped regime. Oscillatory systems require an extended model with delayed feedback, inhibitory coupling, or another state variable.
  \item \textbf{Domain calibration.} The model does not justify precise biological, economic, physical, or computational numerical predictions without independent parameter estimation. The synthetic demonstration in Sec.~\ref{sec:synthetic} illustrates the protocol; it does not validate any real system.
  \item \textbf{Observability.} Many systems do not directly expose both $I(t)$ and $S(t)$. Practical inference may require proxies, latent-variable estimation, or Bayesian fitting.
  \item \textbf{Noise.} Stochastic forcing may shift apparent thresholds, induce switching, or blur the active transition near $R=1$.
\end{enumerate}

\section{Conclusion}
The Adaptive Boundary Equation proposes a minimal active-boundary law for coupled information and entropy dynamics. From Eqs.~\eqref{eq:I}--\eqref{eq:S} follows a clear viability threshold: $\eta>\xi\lambda/\gamma$, equivalently $R=\gamma\eta/(\xi\lambda)>1$. Below threshold, unforced boundary activity is not self-maintaining. Above threshold, a stable positive active branch exists. The active branch is stable whenever it exists. In the minimal positive-coupling form, relaxation is non-oscillatory; oscillatory domains require an extended ABE. The predicted threshold transition and amplitude scaling are visible in synthetic operational telemetry, the critical-slowing trend is verified numerically, and explicit falsification criteria are provided.

The model matters if it remains disciplined. Its value is not that it explains everything. Its value is that it gives adaptive-boundary theory a compact mathematical target that can be simulated, calibrated, falsified, and extended.

\section*{Postscript: a long-range direction}
A natural long-range direction is horizon-like information--entropy relaxation, where the active-boundary interpretation becomes literal rather than metaphorical. The shared structural feature is a boundary that retains information while exchanging entropy with its environment under regulatory constraint, a description that has a formal analogy in black-hole horizon physics, while requiring an entirely separate gravitational derivation against the relevant perturbation theory and lying outside the scope of the present phenomenological treatment. We note the structural resonance only to flag it as a research direction worth separate inquiry.

\section*{Acknowledgments}
The author thanks colleagues and advisors, including AI-assisted technical review, for discussion and criticism during manuscript development. All analytical content, mathematical derivations, and final editorial decisions are the author's own.

\section*{Data and code availability}
A reproducibility archive accompanying this paper provides all scripts required to reproduce Figs.~\ref{fig:phase}--\ref{fig:critical-slowing} and to verify the equilibrium, stability, active-branch, amplitude-scaling, and critical-slowing claims numerically, together with the LaTeX source. The archive is available at the \href{https://github.com/hellocripsis/adaptive-boundary-equation}{project GitHub repository}; an archived snapshot is registered under DOI \href{https://doi.org/10.5281/zenodo.20090457}{10.5281/zenodo.20090457}.

\section*{Competing interests}
The author is affiliated with Cripsis Logic, which is developing commercial systems related to adaptive computational control and information-governance architectures.

\clearpage
\appendix
\section{Complete active-branch derivation}
Starting from the unforced normalized ABE,
\begin{equation*}
\frac{dI}{dt}=\gamma S-\lambda I-\beta I^3,\qquad \frac{dS}{dt}=\eta I-\xi S,
\end{equation*}
at equilibrium $dS/dt=0$ implies $S^*=(\eta/\xi)I^*$. Substitution into the $I$ equation yields
\begin{equation*}
0=(\gamma\eta/\xi-\lambda)I^*-\beta(I^*)^3.
\end{equation*}
Therefore $I^*=0$ or $(I^*)^2=(\gamma\eta/\xi-\lambda)/\beta$. The positive branch exists only if $\gamma\eta/\xi>\lambda$, equivalent to $\eta>\xi\lambda/\gamma$.

The Jacobian of the unforced system is
\begin{equation*}
J(I,S)=\begin{pmatrix}-\lambda-3\beta I^2 & \gamma\\ \eta & -\xi\end{pmatrix}.
\end{equation*}
At the quiescent state, this becomes Eq.~\eqref{eq:J0}. At the active branch, $\beta(I^*)^2=\gamma\eta/\xi-\lambda$, so
\begin{equation*}
-\lambda-3\beta(I^*)^2=2\lambda-3\gamma\eta/\xi,
\end{equation*}
which gives Eq.~\eqref{eq:Jstar}. The determinant is
\begin{align*}
D^*&=(2\lambda-3\gamma\eta/\xi)(-\xi)-\gamma\eta\\
&=-2\lambda\xi+3\gamma\eta-\gamma\eta\\
&=2(\gamma\eta-\lambda\xi).
\end{align*}
For $\eta>\xi\lambda/\gamma$, $D^*>0$. The trace satisfies $T^*=2\lambda-3\gamma\eta/\xi-\xi$. Since $\eta>\xi\lambda/\gamma$ implies $\gamma\eta/\xi>\lambda$,
\begin{equation*}
T^*<2\lambda-3\lambda-\xi=-\lambda-\xi<0.
\end{equation*}
Therefore the positive active branch is stable whenever it exists.

\clearpage
\section{Minimal simulation pseudocode}
\begin{verbatim}
for each parameter set gamma, lambda, beta, eta, xi > 0:
    eta_c = xi * lambda / gamma
    R = eta / eta_c

    if R <= 1:
        assert no positive active branch exists
        J0 = [[-lambda, gamma], [eta, -xi]]
        assert trace(J0) < 0 and det(J0) > 0
    else:
        I_star = sqrt((gamma*eta/xi - lambda) / beta)
        S_star = (eta / xi) * I_star
        J_star = [[(2*lambda - 3*gamma*eta/xi), gamma],
                  [eta, -xi]]
        T = trace(J_star)
        D = det(J_star)
        eig = eigenvalues(J_star)
        assert D > 0 and T < 0
        assert real_part(eig[0]) < 0
        assert real_part(eig[1]) < 0
\end{verbatim}
This pseudocode prevents the most serious submission error: using one set of equations in the model and a different Jacobian in the proof.

\clearpage
\begin{thebibliography}{99}
\bibitem{schrodinger1944life} E. Schr{\"o}dinger, \textit{What is Life?} (Cambridge University Press, Cambridge, 1944).
\bibitem{prigogine1967thermo} I. Prigogine, \textit{Introduction to Thermodynamics of Irreversible Processes} (Wiley, New York, 1967).
\bibitem{nicolis1977self} G. Nicolis and I. Prigogine, \textit{Self-Organization in Nonequilibrium Systems} (Wiley, New York, 1977).
\bibitem{dewar2003info} R.~C. Dewar, J. Phys. A: Math. Gen. \textbf{36}, 631 (2003).
\bibitem{martyushev2006mep} L.~M. Martyushev and V.~D. Seleznev, Phys. Rep. \textbf{426}, 1 (2006).
\bibitem{friston2010free} K. Friston, Nat. Rev. Neurosci. \textbf{11}, 127 (2010).
\bibitem{buckley2017free} C.~L. Buckley, C.~S. Kim, S. McGregor, and A.~K. Seth, J. Math. Psychol. \textbf{81}, 55 (2017).
\bibitem{hawkes1971spectra} A.~G. Hawkes, Biometrika \textbf{58}, 83 (1971).
\bibitem{norris1997markov} J.~R. Norris, \textit{Markov Chains} (Cambridge University Press, Cambridge, 1997).
\bibitem{gardiner2009stochastic} C.~W. Gardiner, \textit{Stochastic Methods: A Handbook for the Natural and Social Sciences} (Springer, Berlin, 2009).
\bibitem{bremaud1996stability} P. Br{\'e}maud and L. Massouli{\'e}, Ann. Probab. \textbf{24}, 1563 (1996).
\bibitem{daley2003point} D.~J. Daley and D. Vere-Jones, \textit{An Introduction to the Theory of Point Processes, Volume I} (Springer, New York, 2003).
\bibitem{fitzhugh1961impulses} R. FitzHugh, Biophys. J. \textbf{1}, 445 (1961).
\bibitem{nagumo1962active} J. Nagumo, S. Arimoto, and S. Yoshizawa, Proc. IRE \textbf{50}, 2061 (1962).
\bibitem{cross1993pattern} M.~C. Cross and P.~C. Hohenberg, Rev. Mod. Phys. \textbf{65}, 851 (1993).
\bibitem{kuramoto1984chemical} Y. Kuramoto, \textit{Chemical Oscillations, Waves, and Turbulence} (Springer, Berlin, 1984).
\bibitem{strogatz2018nonlinear} S.~H. Strogatz, \textit{Nonlinear Dynamics and Chaos}, 2nd ed. (CRC Press, Boca Raton, 2018).
\bibitem{friston2013life} K. Friston, J. R. Soc. Interface \textbf{10}, 20130475 (2013).
\bibitem{kirchhoff2018markov} M.~D. Kirchhoff, T. Parr, E. Palacios, K. Friston, and J. Kiverstein, J. R. Soc. Interface \textbf{15}, 20170792 (2018).
\bibitem{aubin1991viability} J.-P. Aubin, \textit{Viability Theory} (Birkh{\"a}user, Boston, 1991).
\bibitem{transtrum2015perspective} M.~K. Transtrum, B.~B. Machta, K.~S. Brown, B.~C. Daniels, C.~R. Myers, and J.~P. Sethna, J. Chem. Phys. \textbf{143}, 010901 (2015).
\bibitem{daniels2015automated} B.~C. Daniels and I. Nemenman, Nat. Commun. \textbf{6}, 8133 (2015).
\end{thebibliography}

\end{document}